%% ma: L12-B08-0050
%% lop: 12
%% chuong: 2
%% bai: 8
%% dang: 12.8.35
%% loai: DS
%% muc: [NB, NB, TH, VD]
%% dap_an: "ĐĐĐĐ"
%% nguon: "LaTeX giáo viên gửi 21/10/2026 (M05, Vectơ và hệ trục tọa độ trong không gian, 107 câu), Phần II Câu 2"
%% kien_thuc: "Bài 7 (tọa độ điểm), Bài 8 (trung điểm, góc giữa hai vectơ)"
%% trang_thai: chinh-thuc
%% ghi_chu: "Giáo viên đã duyệt 21/10/2026. Hình gốc vẽ lệch trục (nhãn y đặt ở trục sâu, nét khuất sai); đã vẽ lại theo đúng tỉ lệ AB = 1, AD = 2, SA = 3"
\begin{ex}
Cho hình chóp $S.ABCD$ có đáy $ABCD$ là hình chữ nhật, $AB=1$, $AD=2$, cạnh bên $SA$ vuông góc với mặt đáy và $SA=3$. Thiết lập hệ trục tọa độ $Oxyz$ sao cho gốc $O$ trùng với $A$, các tia $Ox$, $Oy$, $Oz$ lần lượt chứa các tia $AB$, $AD$, $AS$ (tham khảo hình bên dưới).
\begin{center}
\begin{tikzpicture}[x={(-1.1cm,-.8cm)},y={(1.3cm,0cm)},z={(0cm,1cm)}]
  \coordinate (A) at (0,0,0); \coordinate (B) at (1,0,0); \coordinate (C) at (1,2,0); \coordinate (D) at (0,2,0); \coordinate (S) at (0,0,3);
  \draw[khuat] (S)--(A)--(B) (A)--(D);
  \draw (S)--(B)--(C)--(D)--cycle (S)--(C);
  \draw[truc] (B)--(1.8,0,0) node[below left]{$x$};
  \draw[truc] (D)--(0,2.8,0) node[right]{$y$};
  \draw[truc] (S)--(0,0,3.7) node[left]{$z$};
  \path (A) node[above left]{$A\equiv O$} (B) node[left]{$B$} (C) node[below right]{$C$} (D) node[above right]{$D$} (S) node[left]{$S$};
\end{tikzpicture}
\end{center}
\choiceTF
{\True Tọa độ đỉnh $B$ là $(1;0;0)$ và đỉnh $S$ là $(0;0;3)$.}
{\True Tọa độ đỉnh $C$ là $(1;2;0)$.}
{\True Gọi $M$ là trung điểm của $SC$. Tọa độ điểm $M$ là $\left(\dfrac12;1;\dfrac32\right)$.}
{\True Côsin của góc tạo bởi hai vectơ $\overrightarrow{SB}$ và $\overrightarrow{SD}$ bằng $\dfrac9{\sqrt{130}}$.}
\loigiai{a) $B(1;0;0)$, $S(0;0;3)$. Đúng.

b) $\overrightarrow{AC}=\overrightarrow{AB}+\overrightarrow{AD}$ nên $C(1;2;0)$. Đúng.

c) $M\left(\dfrac{1+0}2;\dfrac{2+0}2;\dfrac{0+3}2\right)=\left(\dfrac12;1;\dfrac32\right)$. Đúng.

d) $\overrightarrow{SB}=(1;0;-3)$, $\overrightarrow{SD}=(0;2;-3)$. $\cos(\overrightarrow{SB},\overrightarrow{SD})=\dfrac{9}{\sqrt{10}\cdot\sqrt{13}}=\dfrac9{\sqrt{130}}$. Đúng.}
\end{ex}
